% ======================================================================
% Ontology of Continua -- Core 1.3
% Cross-K module: crossk_k11_k12.tex
% Cross-level relations between K11 and a minimally specified K12
% ======================================================================

\subsubsection{Межуровневая структура K11–K12}
\label{sec:crossk-k11-k12}

Переход $K_{11} \rightarrow K_{12}$ является завершающим межуровневым шагом
в Core~1.3. Если $K_{11}$ даёт транс-метатеоретический континуум с
универсальными операторами, действующими на целые классы мета-рамок, то $K_{12}$
объединяет эти верхнеуровневые структуры в финальную оболочку единой науки,
в которой оценивается междоменная совместимость и глобальная интегрируемость.

Этот уровень остаётся абстрактным, но больше не является редакционной заглушкой.
Его роль — зафиксировать условия допустимости, при которых продвигаемая
иерархия может рассматриваться как единый согласованный научный стек.

% ----------------------------------------------------------------------
\subsubsection{1. Задействованные уровни и их роли}

\paragraph{K11 (транс-метатеоретический континуум).}
$K_{11}$ предоставляет:
\begin{itemize}
  \item оси $A^{11}$, относящиеся к классам мета-рамок,
  \item потенциалы $P^{11}$, выражающие транс-мета согласованность,
  \item пороги $\Theta^{11}$ для универсальной согласованности и совместимости,
  \item потоки $J^{11}$, перестраивающие семейства объектов $K_{11}$,
  \item циклы $C^{11}$, регулирующие эволюцию структур на мета-уровне.
\end{itemize}

\paragraph{K12 (формальный верхний континуум).}
В Core~1.3 $K_{12}$ задаётся через:
\begin{itemize}
  \item существование окружающего пространства $M_{12}$ возможных осей,
  \item возможность новых различий, не представимых в $A^{11}$,
  \item абстрактные пороги $\Theta^{12}$, ограничивающие структурную
        согласованность,
  \item обобщённые потоки $J^{12}$ на классах объектов $K_{11}$,
  \item потенциальные циклы $C^{12}$, стабилизирующие новые высшие структуры.
\end{itemize}

$K_{12}$ не является конкретной предметной моделью. Это глобальный слой
интегрируемости, который ограничивает, как замкнутая иерархия может быть
объединена без противоречий.

% ----------------------------------------------------------------------
\subsubsection{2. Наследованные / совместные оси и пороги}

\paragraph{Наследование.}
Все оси, потенциалы и пороги $K_{11}$ вложимы в пространстве
$K_{12}$:
\[
A^{11} \subseteq M_{12}, \qquad
\Theta^{11} \subseteq \Theta^{12}.
\]

\paragraph{Новые различия.}
Если $K_{12}$ существует, его оси должны удовлетворять:
\[
A^{12} \in M_{12} \setminus A(K_{11}),
\]
в соответствии с Теоремой~4 (новые оси не могут быть самообразованными в $K_{11}$).

\paragraph{Расширение порогов.}
$\Theta^{12}$ может включать:
\begin{itemize}
  \item \textbf{расширенные пороги когерентности}, охватывающие целые цепочки
        $K_0\dots K_{11}$,
  \item \textbf{верхние ограничения} (лимиты, налагаемые $M_{12}$),
  \item \textbf{пороги совместимости} для структур выше мета-уровня.
\end{itemize}

% ----------------------------------------------------------------------
\subsubsection{3. Межуровневые потоки, циклы и напряжения}

\paragraph{Потоки.}
В абстрактной форме достаточно:
\[
J^{12} = (J_{\mathrm{lift3}},\ J_{\mathrm{ultra2}},\
          J_{\mathrm{constraint2}},\ J_{\mathrm{embed}}),
\]
что означает: $K_{12}$ действует на структуру $K_{11}$ аналогично тому, как
$K_{11}$ действует на $K_{10}$, но на более высоком уровне.

\paragraph{Циклы.}
Потенциальные циклы $C^{12}$ обобщают транс-мета-циклы:
\[
C_{\mathrm{limit}}:\
\text{K11-structure} \rightarrow \text{lift} \rightarrow
\text{constraint} \rightarrow \text{refinement} \rightarrow \text{K11-structure}.
\]

\paragraph{Напряжение.}
Межуровневое напряжение следует обычной форме:
\[
T_{11,12} =
f(\Theta^{12} - P^{12},\ A^{12} - A^{11},\
  J^{12} - J^{11},\ \Theta^{11} - \Theta^{12}).
\]
Избыточное напряжение коллапсирует $K_{12}$ в $K_{11}$.

% ----------------------------------------------------------------------
\subsubsection{4. Условия рождения и исчезновения между уровнями}

\paragraph{Рождение $K_{12}$.}
В формализованной постановке:
\[
T_{11} > \Theta_{11,\mathrm{dim}}
\quad\text{и}\quad
A^{12} \in M_{12},
\]
то есть в $M_{12}$ должны существовать новые различения, превосходящие
представительностные возможности $K_{11}$.

Расширение пространства состояний:
\[
\Omega(K_{12}) = \Omega(K_{11}) \cup \Delta\Omega_{\mathrm{limit}}.
\]

\paragraph{Распространение гибели.}
Если $\Theta^{11}$ нарушена, структура $K_{12}$ не может быть сформирована.

Если нарушаются $\Theta^{12}$ (расширенная несовместимость или парадокс), тогда:
\[
\Omega(K_{12}) \rightarrow \varnothing,
\]
в то время как $K_{11}$ остаётся целостным.

% ----------------------------------------------------------------------
\subsubsection{5. Концептуальные примеры}

\paragraph{Расширение мета-предела.}
$K_{12}$ может выражать свойства, относящиеся ко всем цепочкам отображений
по уровням $K_0\dots K_{11}$.

\paragraph{Ультра-согласованность.}
Универсальное ограничение может требовать согласованности всей модели под
преобразованиями, не представимыми в $K_{11}$.

\paragraph{Случай коллапса.}
Любое нарушение расширенной когерентности:
\[
T_{11,12} > \Theta^{12}_{\mathrm{coh}},
\]
коллапсирует гипотетический $K_{12}$ обратно в полностью специфицированный
$K_{11}$.

Следовательно, в пределах Core~1.3 переход K11→K12 можно описать в сжатой форме
как формальный верхний контур для междоменной совместимости и глобальной
интегрируемости, гарантируя, что продвигаемая иерархия остаётся научно
согласованной при соблюдении непрерывности и аксиом размерности.
